\section{Conclusion}
\label{sec:conclusion}

We began with an observation that should give the LLM evaluation community pause: a program that outputs the correct answer on 764 tests can still be wrong 40\% of the time --- if you ask a formal contract on the inputs that matter most to contract semantics. This paper has provided the rigorous experimental support for that claim: across 364 ContractEval tasks, 3 model families (Experiment A), and 10,432 evaluation triples, the contract oracle detects 40\% more failures than the differential oracle on contract-violating inputs, with 40\% contract-unique mass. Adaptive PBT results span all 5 tested model families (Experiment B; 17,226 triples).

Our main contributions are: (1) the oracle isolation design --- the first methodology for clean oracle-type comparison via CVT inputs; (2) the 40\% oracle-isolation gap (95\% CI [0.358, 0.445], $p = 5.88 \times 10^{-38}$) confirmed across all tasks and models; (3) the model-consistent $\sim$10pp adaptive PBT contribution (uniform across all 5 tested families) ($p = 2.64 \times 10^{-22}$); and (4) principled negative results on contract richness gradient and cross-model orthogonality.

Future directions include expanding to $n \geq 10$ models for adequate $\tau$ statistical power, filtering postcondition-only contracts for the richness gradient analysis, and testing whether the oracle gap persists on random (non-CVT) input distributions.

Returning to our opening: the gap between ``passes dense differential testing'' and ``satisfies formal contracts'' is structural, not marginal --- 40\% on the inputs most relevant for contract evaluation. As LLM coding capabilities advance, contract-based evaluation offers a principled path to specification-level assessment that differential testing cannot replicate. We hope this work encourages oracle-type diversity as a first-class dimension in LLM code evaluation benchmarks.
